%% Worksheets, self checks and unit test, AP Physics 1 Unit 1 (Kinematics), matching the guided notes 1.1 to 1.5.
%% Use g = 10 m/s^2 and ignore air resistance unless stated.
\renewcommand{\nocalc}{}
\newcommand{\gr}[2]{\begin{center}\begin{tikzpicture}\begin{axis}[gridaxis,#1]#2\end{axis}\end{tikzpicture}\end{center}}
\newcommand{\pt}[1]{\node[circle,fill=apblue,inner sep=1.5pt] at (axis cs:#1){};}
\newcommand{\work}[1]{\ifkey\else\par\vspace{#1}\fi}
\newcommand{\gfig}[1]{\Needspace{20\baselineskip}\begin{center}\textsf{\small\bfseries\color{apnavy}#1}\end{center}}
\newcommand{\ms}{\,\mathrm{m/s}}
\newcommand{\mss}{\,\mathrm{m/s^2}}
\newcommand{\m}{\,\mathrm{m}}
\newcommand{\s}{\,\mathrm{s}}
% 1.3: position-time graph
\newcommand{\gXa}{\gr{xmin=-0.3,xmax=8.6,ymin=-4.5,ymax=7.5,width=8.4cm,height=5.4cm,xlabel={$t$ (s)},ylabel={$x$ (m)},xtick={0,1,...,8},ytick={-3,0,3,6}}{\addplot[apblue,very thick] coordinates {(0,0)(2,6)(5,6)(8,-3)};}}
% 1.3: velocity-time graph
\newcommand{\gVa}{\gr{xmin=-0.3,xmax=12.6,ymin=-5,ymax=9.5,width=9cm,height=5.6cm,xlabel={$t$ (s)},ylabel={$v$ (m/s)},xtick={0,2,...,12},ytick={-4,0,4,8}}{\addplot[apblue,very thick] coordinates {(0,0)(4,8)(6,8)(12,-4)};}}
% pretest: velocity-time graph
\newcommand{\gVb}{\gr{xmin=-0.3,xmax=8.6,ymin=-5,ymax=5.5,width=8.4cm,height=5.2cm,xlabel={$t$ (s)},ylabel={$v$ (m/s)},xtick={0,1,...,8},ytick={-4,-2,0,2,4}}{\addplot[apblue,very thick] coordinates {(0,-4)(4,4)(8,4)};}}
% blank axes for sketches
\newcommand{\blank}[2]{\gr{xmin=0,xmax=#1,ymin=0,ymax=1,width=7cm,height=4cm,xlabel={$t$ (s)},ylabel={#2},xticklabels={},yticklabels={},grid=none}{}}

%%DOC 1.1A | Scalars and Vectors in One Dimension
\begin{exo}{}
Which of the following quantities is a vector?
\choices{Speed}{Distance}{Displacement}{Time}
\answer{C}{Displacement has a magnitude and a direction. Speed, distance and time have a magnitude only.}
\end{exo}
\begin{exo}{}
Which of the following quantities is a scalar?
\choices{Velocity}{Acceleration}{Force}{Mass}
\answer{D}{Mass has no direction. Velocity, acceleration and force are vectors.}
\end{exo}
\begin{exo}{}
An object moves along the $x$-axis from $x=+3\m$ to $x=-5\m$. What is its displacement?
\choices{$+8\m$}{$-8\m$}{$-2\m$}{$+2\m$}
\answer{B}{$\Delta x=x_f-x_i=-5-(+3)=-8\m$.}
\end{exo}
\begin{exo}{}
A person walks $30\m$ east and then $50\m$ west. What are the distance traveled and the displacement?
\choices{$80\m$; $20\m$ west}{$80\m$; $20\m$ east}{$20\m$; $80\m$ west}{$20\m$ west; $80\m$}
\answer{A}{Distance $=30+50=80\m$. Displacement $=+30-50=-20\m$, that is $20\m$ west.}
\end{exo}
\begin{exo}{}
East is positive. Which velocity has the greatest speed?
\choices{$+6\ms$}{$-9\ms$}{$+8\ms$}{$-4\ms$}
\answer{B}{Speed is the magnitude: $9\ms$ is the largest. The minus sign only means the motion is to the west.}
\end{exo}
\begin{exo}{}
A ball is thrown straight up and caught at the same height. What is its displacement for the whole trip?
\choices{Twice the maximum height}{Zero}{The maximum height}{Negative}
\answer{B}{It ends where it started, so $\Delta y=0$, even though the distance traveled is twice the maximum height.}
\end{exo}
\begin{exo}{}
Up is positive. A stone falls $12\m$ from the top of a cliff to the ground. What is its displacement?
\choices{$+12\m$}{$0\m$}{$-12\m$}{$24\m$}
\answer{C}{It moves $12\m$ in the negative (downward) direction.}
\end{exo}
\begin{exo}{}
Two forces act on a box along the same line: $40\,\mathrm{N}$ to the right and $25\,\mathrm{N}$ to the left. What is the net force?
\choices{$65\,\mathrm{N}$ to the right}{$15\,\mathrm{N}$ to the left}{$15\,\mathrm{N}$ to the right}{$65\,\mathrm{N}$ to the left}
\answer{C}{$+40-25=+15\,\mathrm{N}$: to the right.}
\end{exo}
\Needspace{12\baselineskip}
\begin{exo}[Free response]{}
A car on a straight road starts at $x=0$, drives to $x=+120\m$, then back to $x=+40\m$, then forward to $x=+70\m$.
\begin{enumerate}[label=(\alph*)]
\item Find the total distance traveled.
\item Find the displacement.
\item Can the magnitude of a displacement ever be larger than the distance traveled? Explain.
\end{enumerate}
\work{4cm}
\begin{solution}
(a) $120+80+30=230\m$. \;(b) $\Delta x=70-0=+70\m$. \;(c) No. The displacement is the straight-line change in position; the path is at least that long. They are equal only if the object never turns around.
\end{solution}
\end{exo}

%%DOC 1.2A | Displacement, Velocity, and Acceleration
\begin{exo}{}
A runner covers $100\m$ in a straight line in $12.5\s$. What is her average speed?
\choices{$12.5\ms$}{$80\ms$}{$0.125\ms$}{$8\ms$}
\answer{D}{$100/12.5=8\ms$.}
\end{exo}
\begin{exo}{}
A car speeds up from $10\ms$ to $25\ms$ in $5\s$. What is its average acceleration?
\choices{$2\mss$}{$5\mss$}{$3\mss$}{$7\mss$}
\answer{C}{$a=\Delta v/\Delta t=(25-10)/5=3\mss$.}
\end{exo}
\begin{exo}{}
An object has velocity $v=-6\ms$ and acceleration $a=+2\mss$. What is happening?
\choices{It is speeding up in the negative direction.}{It is slowing down.}{It is at rest.}{It is speeding up in the positive direction.}
\answer{B}{$v$ and $a$ have opposite signs, so the object slows down.}
\end{exo}
\begin{exo}{}
An object starts from rest with a constant acceleration of $3\mss$. What is its speed after $4\s$?
\choices{$7\ms$}{$12\ms$}{$24\ms$}{$0.75\ms$}
\answer{B}{$v=v_0+at=0+3(4)=12\ms$.}
\end{exo}
\begin{exo}{}
For the object in the previous question, how far does it travel in those $4\s$?
\choices{$12\m$}{$48\m$}{$24\m$}{$6\m$}
\answer{C}{$\Delta x=\frac12at^2=\frac12(3)(16)=24\m$.}
\end{exo}
\begin{exo}{}
A car moving at $30\ms$ brakes with an acceleration of magnitude $6\mss$. How far does it travel before stopping?
\choices{$75\m$}{$150\m$}{$5\m$}{$90\m$}
\answer{A}{$0=30^2-2(6)\Delta x$, so $\Delta x=900/12=75\m$.}
\end{exo}
\begin{exo}{}
A ball is thrown straight up at $15\ms$. How long does it take to reach its highest point?
\choices{$3\s$}{$0.67\s$}{$1.5\s$}{$15\s$}
\answer{C}{At the top $v=0$: $0=15-10t$, so $t=1.5\s$.}
\end{exo}
\begin{exo}{}
At the highest point of its path, a ball thrown straight up has
\choices{$v=0$ and $a=0$}{$v=0$ and $a=10\mss$ downward}{$v=10\ms$ upward and $a=0$}{$v=0$ and $a=10\mss$ upward}
\answer{B}{The velocity is zero for an instant, but gravity still acts: $a=g$ downward at every point.}
\end{exo}
\Needspace{12\baselineskip}
\begin{exo}[Free response]{}
A train starts from rest and accelerates at $0.50\mss$ for $40\s$. It then brakes with an acceleration of magnitude $1.0\mss$ until it stops.
\begin{enumerate}[label=(\alph*)]
\item Find its speed at $t=40\s$.
\item Find the distance traveled during the first $40\s$.
\item Find the braking distance and the braking time.
\end{enumerate}
\work{4.5cm}
\begin{solution}
(a) $v=0.5(40)=20\ms$. \;(b) $\Delta x=\frac12(0.5)(40)^2=400\m$. \;(c) $0=20^2-2(1.0)\Delta x\Rightarrow\Delta x=200\m$; $t=20/1.0=20\s$.
\end{solution}
\end{exo}

%%DOC 1.3A | Representing Motion
\gfig{Questions 1 to 4 refer to the graph below: the position-time graph of an object moving along a line.}
\gXa
\begin{exo}{}
What is the velocity of the object between $t=0$ and $t=2\s$?
\choices{$6\ms$}{$3\ms$}{$2\ms$}{$0$}
\answer{B}{Slope $=6/2=3\ms$.}
\end{exo}
\begin{exo}{}
Between $t=2\s$ and $t=5\s$, the object is
\choices{at rest}{moving at a constant positive velocity}{speeding up}{moving backward}
\answer{A}{The position stays at $6\m$: the slope is zero, so the velocity is zero.}
\end{exo}
\begin{exo}{}
What is the velocity of the object between $t=5\s$ and $t=8\s$?
\choices{$3\ms$}{$-9\ms$}{$-3\ms$}{$1\ms$}
\answer{C}{Slope $=(-3-6)/(8-5)=-3\ms$: it moves in the negative direction.}
\end{exo}
\begin{exo}{}
What is the displacement of the object from $t=0$ to $t=8\s$?
\choices{$6\m$}{$-3\m$}{$15\m$}{$3\m$}
\answer{B}{$\Delta x=x(8)-x(0)=-3-0=-3\m$.}
\end{exo}
\gfig{Questions 5 to 8 refer to the graph below: the velocity-time graph of a cart moving along a line.}
\gVa
\begin{exo}{}
What is the acceleration of the cart at $t=2\s$?
\choices{$8\mss$}{$4\mss$}{$2\mss$}{$0$}
\answer{C}{Slope from $0$ to $4\s$: $8/4=2\mss$.}
\end{exo}
\begin{exo}{}
What is the displacement of the cart from $t=0$ to $t=6\s$?
\choices{$32\m$}{$48\m$}{$16\m$}{$24\m$}
\answer{A}{Area: triangle $\frac12(4)(8)=16\m$ plus rectangle $2\times8=16\m$: $32\m$.}
\end{exo}
\begin{exo}{}
At what time does the cart change direction?
\choices{$6\s$}{$4\s$}{$12\s$}{$10\s$}
\answer{D}{The velocity changes sign where the graph crosses the $t$-axis: at $t=10\s$.}
\end{exo}
\begin{exo}{}
What is the displacement of the cart from $t=0$ to $t=12\s$?
\choices{$52\m$}{$48\m$}{$44\m$}{$40\m$}
\answer{C}{Areas: $16+16+16$ above the axis (to $t=10\s$), then $-\frac12(2)(4)=-4\m$ below: $48-4=44\m$. (The distance is $52\m$.)}
\end{exo}
\Needspace{24\baselineskip}
\begin{exo}[Free response]{}
A cart starts from rest at $x=0$. It accelerates at $2\mss$ for $3\s$, then moves at constant velocity for $3\s$.
\begin{enumerate}[label=(\alph*)]
\item On the first axes, sketch the velocity-time graph from $t=0$ to $t=6\s$. Label the values.
\item On the second axes, sketch the position-time graph. Describe the shape of each part.
\item Find the displacement at $t=6\s$.
\end{enumerate}
\begin{minipage}{0.48\linewidth}\blank{6}{$v$}\end{minipage}\hfill\begin{minipage}{0.48\linewidth}\blank{6}{$x$}\end{minipage}
\work{2.5cm}
\begin{solution}
(a) A straight line from $(0,0)$ to $(3\s,6\ms)$, then horizontal at $6\ms$. \;(b) A parabola opening up (slope increasing) from $0$ to $3\s$, reaching $9\m$, then a straight line of slope $6\ms$. \;(c) $\frac12(3)(6)+3(6)=9+18=27\m$.
\end{solution}
\end{exo}

%%DOC 1.4A | Reference Frames and Relative Motion
\begin{exo}{}
A train moves at $25\ms$. A passenger walks toward the front of the train at $2\ms$ relative to the train. What is the passenger's speed relative to the ground?
\choices{$23\ms$}{$27\ms$}{$25\ms$}{$2\ms$}
\answer{B}{Same direction: $25+2=27\ms$.}
\end{exo}
\begin{exo}{}
The same passenger now walks toward the back of the train at $2\ms$ relative to the train. Her velocity relative to the ground is
\choices{$23\ms$ forward}{$27\ms$ backward}{$23\ms$ backward}{$2\ms$ backward}
\answer{A}{$+25-2=+23\ms$: she still moves forward relative to the ground.}
\end{exo}
\begin{exo}{}
Car A moves north at $30\ms$ and car B moves north at $20\ms$. What is the velocity of B relative to A?
\choices{$10\ms$ north}{$50\ms$ north}{$10\ms$ south}{$50\ms$ south}
\answer{C}{$\vec v_{B/A}=\vec v_B-\vec v_A=20-30=-10\ms$: seen from A, car B moves backward (south).}
\end{exo}
\begin{exo}{}
Two cars drive toward each other, each at $20\ms$ relative to the road. What is the speed of one car relative to the other?
\choices{$0$}{$20\ms$}{$10\ms$}{$40\ms$}
\answer{D}{Opposite directions: $20-(-20)=40\ms$.}
\end{exo}
\begin{exo}{}
A boat points straight across a river $60\m$ wide at $5\ms$ relative to the water. The current is $2\ms$. How long does it take to cross?
\choices{$12\s$}{$30\s$}{$8.6\s$}{$11.1\s$}
\answer{A}{Only the velocity across the river matters: $60/5=12\s$.}
\end{exo}
\begin{exo}{}
How far downstream does the boat in the previous question land?
\choices{$12\m$}{$24\m$}{$60\m$}{$30\m$}
\answer{B}{The current carries it $2\ms\times12\s=24\m$ downstream.}
\end{exo}
\begin{exo}{}
A passenger in a train moving at constant velocity drops a ball. For the passenger, the path of the ball is
\choices{straight down}{a curve bending backward}{a curve bending forward}{horizontal}
\answer{A}{In the train's frame the ball starts at rest and only falls: straight down.}
\end{exo}
\begin{exo}{}
For an observer standing beside the track, the path of the same ball is
\choices{straight down}{a curve that moves forward with the train}{a curve that moves backward}{a horizontal line}
\answer{B}{In the ground frame the ball also has the train's horizontal velocity: it follows a projectile path forward.}
\end{exo}
\Needspace{12\baselineskip}
\begin{exo}[Free response]{}
An airplane heads north at $200\ms$ relative to the air. A wind blows east at $50\ms$.
\begin{enumerate}[label=(\alph*)]
\item Find the magnitude and direction of the plane's velocity relative to the ground.
\item How far north does the plane travel in one hour?
\item In which direction should the pilot point the plane to fly due north? Explain.
\end{enumerate}
\work{4.5cm}
\begin{solution}
(a) $\sqrt{200^2+50^2}\approx206\ms$, at $\tan^{-1}(50/200)\approx14^\circ$ east of north. \;(b) $200\times3600=720\,000\m=720\,\mathrm{km}$. \;(c) Slightly west of north, so that the westward component of the plane's velocity cancels the wind.
\end{solution}
\end{exo}

%%DOC 1.5A | Vectors and Motion in Two Dimensions
\begin{exo}{}
A ball is launched at $25\ms$ at $37^\circ$ above the horizontal ($\cos37^\circ=0.8$, $\sin37^\circ=0.6$). What is the horizontal component of its velocity?
\choices{$15\ms$}{$20\ms$}{$25\ms$}{$12.5\ms$}
\answer{B}{$v_x=25\cos37^\circ=25(0.8)=20\ms$.}
\end{exo}
\begin{exo}{}
A vector has components $A_x=6$ and $A_y=-8$. What is its magnitude?
\choices{$14$}{$2$}{$10$}{$48$}
\answer{C}{$\sqrt{6^2+8^2}=10$.}
\end{exo}
\begin{exo}{}
A ball is thrown horizontally at $8\ms$ from the top of a cliff and lands $2\s$ later. How high is the cliff?
\choices{$20\m$}{$16\m$}{$10\m$}{$40\m$}
\answer{A}{Vertical: $h=\frac12gt^2=\frac12(10)(4)=20\m$. The horizontal speed does not matter.}
\end{exo}
\begin{exo}{}
How far from the base of the cliff does the ball in the previous question land?
\choices{$8\m$}{$4\m$}{$32\m$}{$16\m$}
\answer{D}{Horizontal: $x=v_xt=8(2)=16\m$.}
\end{exo}
\begin{exo}{}
At the same moment and from the same height, ball 1 is dropped and ball 2 is thrown horizontally. Which lands first?
\choices{Ball 1}{Ball 2}{They land at the same time.}{It depends on their masses.}
\answer{C}{Both start with $v_y=0$ and have the same vertical acceleration $g$, so they fall the same height in the same time.}
\end{exo}
\begin{exo}{}
What is the acceleration of a projectile at the highest point of its path?
\choices{Zero}{$g$ downward}{$g$ upward}{Horizontal}
\answer{B}{Gravity acts during the whole flight: $a=g$ downward, even at the top.}
\end{exo}
\begin{exo}{}
A ball is launched at $30\ms$ at $53^\circ$ above level ground ($\cos53^\circ=0.6$, $\sin53^\circ=0.8$). How long is it in the air?
\choices{$2.4\s$}{$4.8\s$}{$3\s$}{$6\s$}
\answer{B}{$v_{0y}=24\ms$; time to the top $24/10=2.4\s$; on level ground the total time is twice that: $4.8\s$.}
\end{exo}
\begin{exo}{}
Projectiles are launched with the same speed from level ground. Which launch angle gives the greatest range?
\choices{$30^\circ$}{$60^\circ$}{$45^\circ$}{$90^\circ$}
\answer{C}{Without air resistance the range is greatest at $45^\circ$ ($30^\circ$ and $60^\circ$ give the same, smaller range).}
\end{exo}
\Needspace{12\baselineskip}
\begin{exo}[Free response]{}
A soccer ball is kicked from the ground at $20\ms$ at $37^\circ$ above the horizontal ($\cos37^\circ=0.8$, $\sin37^\circ=0.6$).
\begin{enumerate}[label=(\alph*)]
\item Find the horizontal and vertical components of the initial velocity.
\item Find the time to reach the highest point and the maximum height.
\item Find the total time in the air and the horizontal range.
\item What is the speed of the ball at its highest point?
\end{enumerate}
\work{5cm}
\begin{solution}
(a) $v_x=16\ms$, $v_{0y}=12\ms$. \;(b) $t=12/10=1.2\s$; $h=\frac{12^2}{2(10)}=7.2\m$. \;(c) $2.4\s$; $16\times2.4=38.4\m$. \;(d) Only $v_x$ remains: $16\ms$.
\end{solution}
\end{exo}

%%DOC SC1.1 | Self Check 1.1 \textbullet\ Topics 1.1 to 1.3
\begin{exo}{}
Which pair contains only vector quantities?
\choices{Displacement and velocity}{Speed and velocity}{Distance and acceleration}{Time and force}
\answer{A}{Displacement and velocity both have a direction. Speed, distance and time are scalars.}
\end{exo}
\begin{exo}{}
A cyclist rides $6\,\mathrm{km}$ east, then $2\,\mathrm{km}$ west, in $0.5\,\mathrm{h}$. What is the average velocity?
\choices{$16\,\mathrm{km/h}$ east}{$8\,\mathrm{km/h}$ east}{$8\,\mathrm{km/h}$ west}{$4\,\mathrm{km/h}$ east}
\answer{B}{Displacement $4\,\mathrm{km}$ east, divided by $0.5\,\mathrm{h}$: $8\,\mathrm{km/h}$ east. (The average speed is $16\,\mathrm{km/h}$.)}
\end{exo}
\begin{exo}{}
A car moving in the negative $x$-direction is slowing down. Which signs are correct?
\choices{$v<0$, $a<0$}{$v>0$, $a<0$}{$v<0$, $a>0$}{$v>0$, $a>0$}
\answer{C}{Moving in the negative direction: $v<0$. Slowing down: $a$ has the opposite sign, $a>0$.}
\end{exo}
\begin{exo}{}
An object starts from rest and accelerates at $4\mss$ over a distance of $18\m$. What is its final speed?
\choices{$6\ms$}{$72\ms$}{$144\ms$}{$12\ms$}
\answer{D}{$v^2=0+2(4)(18)=144$, so $v=12\ms$.}
\end{exo}
\begin{exo}{}
On a velocity-time graph, the area between the graph and the time axis represents the
\choices{acceleration}{displacement}{speed}{time}
\answer{B}{Velocity $\times$ time $=$ displacement.}
\end{exo}
\begin{exo}{}
The position-time graph of an object is a horizontal line. The object is
\choices{at rest}{moving with constant acceleration}{moving at a constant nonzero velocity}{slowing down}
\answer{A}{The slope (velocity) is zero.}
\end{exo}
\Needspace{12\baselineskip}
\begin{exo}[Free response]{}
An elevator starts from rest and moves up. It reaches $3\ms$ in $2\s$, moves at $3\ms$ for $4\s$, then stops in $3\s$.
\begin{enumerate}[label=(\alph*)]
\item Find the acceleration during each part.
\item Find the total displacement.
\item Find the average velocity for the whole trip.
\end{enumerate}
\work{4.5cm}
\begin{solution}
(a) $1.5\mss$; $0$; $-1\mss$. \;(b) $\frac12(2)(3)+4(3)+\frac12(3)(3)=3+12+4.5=19.5\m$ up. \;(c) $19.5/9\approx2.2\ms$ up.
\end{solution}
\end{exo}

%%DOC SC1.2 | Self Check 1.2 \textbullet\ Topics 1.4 and 1.5
\begin{exo}{}
A boat moves straight across a river at $6\ms$ relative to the water. The current is $8\ms$. What is the boat's speed relative to the bank?
\choices{$14\ms$}{$2\ms$}{$10\ms$}{$48\ms$}
\answer{C}{The two velocities are perpendicular: $\sqrt{6^2+8^2}=10\ms$.}
\end{exo}
\begin{exo}{}
Rain falls straight down at $8\ms$. You run at $6\ms$. At what speed does the rain hit you, in your reference frame?
\choices{$2\ms$}{$10\ms$}{$14\ms$}{$8\ms$}
\answer{B}{Relative to you the rain also moves backward at $6\ms$: $\sqrt{8^2+6^2}=10\ms$.}
\end{exo}
\begin{exo}{}
A ball is thrown horizontally from the top of a building $45\m$ tall. How long does it take to reach the ground?
\choices{$3\s$}{$4.5\s$}{$9\s$}{$2\s$}
\answer{A}{$45=\frac12(10)t^2\Rightarrow t^2=9$, $t=3\s$.}
\end{exo}
\begin{exo}{}
During the flight of a projectile (no air resistance), the horizontal component of its velocity
\choices{increases}{decreases}{stays constant}{is zero at the top}
\answer{C}{There is no horizontal acceleration.}
\end{exo}
\begin{exo}{}
A projectile is launched at $40\ms$ at $30^\circ$ above the horizontal. What is its maximum height above the launch point?
\choices{$20\m$}{$80\m$}{$40\m$}{$10\m$}
\answer{A}{$v_{0y}=40\sin30^\circ=20\ms$; $h=\frac{20^2}{2(10)}=20\m$.}
\end{exo}
\begin{exo}{}
A force of $10\,\mathrm{N}$ acts at $37^\circ$ above the horizontal ($\sin37^\circ=0.6$). What is its vertical component?
\choices{$8\,\mathrm{N}$}{$6\,\mathrm{N}$}{$10\,\mathrm{N}$}{$4\,\mathrm{N}$}
\answer{B}{$10\sin37^\circ=6\,\mathrm{N}$.}
\end{exo}
\Needspace{12\baselineskip}
\begin{exo}[Free response]{}
A plane flies horizontally at $60\ms$ at a height of $500\m$ and releases a package.
\begin{enumerate}[label=(\alph*)]
\item How long does the package take to reach the ground?
\item How far does it travel horizontally?
\item Where is the plane when the package lands? Explain using reference frames.
\end{enumerate}
\work{4.5cm}
\begin{solution}
(a) $500=\frac12(10)t^2\Rightarrow t=10\s$. \;(b) $60\times10=600\m$. \;(c) Directly above it: the package keeps the plane's horizontal velocity, so in the plane's frame it falls straight down.
\end{solution}
\end{exo}

%%DOC PRETEST1 | Unit 1 Test \textbullet\ Kinematics
\partheading{Part A \textbullet\ Multiple choice}
\begin{exo}{}
Which quantity has both a magnitude and a direction?
\choices{Speed}{Distance}{Acceleration}{Time}
\answer{C}{Acceleration is a vector.}
\end{exo}
\begin{exo}{}
A ball rolls $4\m$ east, then $7\m$ west. What is its displacement?
\choices{$11\m$ east}{$3\m$ west}{$3\m$ east}{$11\m$ west}
\answer{B}{$+4-7=-3\m$: $3\m$ west.}
\end{exo}
\begin{exo}{}
A car speeds up uniformly from rest to $20\ms$ in $8\s$. What is its acceleration?
\choices{$2.5\mss$}{$160\mss$}{$0.4\mss$}{$12\mss$}
\answer{A}{$20/8=2.5\mss$.}
\end{exo}
\begin{exo}{}
How far does the car in the previous question travel in those $8\s$?
\choices{$160\m$}{$40\m$}{$80\m$}{$20\m$}
\answer{C}{Average velocity $(0+20)/2=10\ms$, times $8\s$: $80\m$.}
\end{exo}
\begin{exo}{}
An object has $v=+5\ms$ and a constant acceleration $a=-2\mss$. What is its velocity $4\s$ later?
\choices{$+13\ms$}{$-3\ms$}{$+3\ms$}{$-13\ms$}
\answer{B}{$v=5+(-2)(4)=-3\ms$: it has turned around.}
\end{exo}
\begin{exo}{}
A rock is dropped from rest and falls for $3\s$. How far does it fall?
\choices{$30\m$}{$90\m$}{$15\m$}{$45\m$}
\answer{D}{$\frac12(10)(3)^2=45\m$.}
\end{exo}
\begin{exo}{}
A ball is thrown straight up at $20\ms$. What maximum height does it reach?
\choices{$40\m$}{$20\m$}{$10\m$}{$2\m$}
\answer{B}{$h=\frac{v_0^2}{2g}=\frac{400}{20}=20\m$.}
\end{exo}
\gfig{Questions 8 to 11 refer to the graph below: the velocity-time graph of an object moving along a line.}
\gVb
\begin{exo}{}
What is the acceleration of the object between $t=0$ and $t=4\s$?
\choices{$0$}{$1\mss$}{$2\mss$}{$4\mss$}
\answer{C}{Slope $=(4-(-4))/4=2\mss$.}
\end{exo}
\begin{exo}{}
At what time is the object momentarily at rest?
\choices{$0$}{$2\s$}{$4\s$}{$8\s$}
\answer{B}{$v=0$ where the graph crosses the axis: $t=2\s$.}
\end{exo}
\begin{exo}{}
What is the displacement from $t=0$ to $t=4\s$?
\choices{$0$}{$8\m$}{$16\m$}{$-4\m$}
\answer{A}{Area below the axis $-4\m$ (from $0$ to $2\s$) plus area above $+4\m$ (from $2$ to $4\s$): $0$. The object comes back to its starting point.}
\end{exo}
\begin{exo}{}
What is the displacement from $t=0$ to $t=8\s$?
\choices{$24\m$}{$32\m$}{$16\m$}{$20\m$}
\answer{C}{$0$ for the first $4\s$, then $4\times4=16\m$.}
\end{exo}
\begin{exo}{}
An escalator moves up at $0.5\ms$. A person walks up the escalator at $1.0\ms$ relative to the steps. What is the person's speed relative to the building?
\choices{$1.5\ms$}{$0.5\ms$}{$1.0\ms$}{$2.0\ms$}
\answer{A}{Same direction: $1.0+0.5=1.5\ms$.}
\end{exo}
\begin{exo}{}
A ball is thrown horizontally at $12\ms$ from a height of $5\m$. How far from the launch point, horizontally, does it land?
\choices{$6\m$}{$24\m$}{$60\m$}{$12\m$}
\answer{D}{$5=\frac12(10)t^2\Rightarrow t=1\s$; $x=12\times1=12\m$.}
\end{exo}
\begin{exo}{}
Which quantity stays constant during the flight of a projectile (no air resistance)?
\choices{The vertical velocity}{The speed}{The horizontal velocity}{The height}
\answer{C}{$a_x=0$, so $v_x$ is constant. The vertical velocity, the speed and the height change.}
\end{exo}
\Needspace{24\baselineskip}
\partheading{Part B \textbullet\ Free response}
\Needspace{16\baselineskip}
\begin{exo}[Free response \textbullet\ Mathematical routines]{}
A ball is thrown straight up at $10\ms$ from the edge of a roof $15\m$ above the ground. It misses the roof on the way down and lands on the ground.
\begin{enumerate}[label=(\alph*)]
\item Find the time to reach the highest point and the maximum height above the ground.
\item Find the total time from the throw until the ball hits the ground.
\item Find the speed of the ball just before it hits the ground.
\item Sketch the velocity-time graph of the ball (up is positive), from the throw to the impact. Label the values.
\end{enumerate}
\blank{3}{$v$}
\work{5cm}
\begin{solution}
(a) $t=10/10=1\s$; it rises $\frac{10^2}{20}=5\m$, so $15+5=20\m$ above the ground. \;(b) From the top it falls $20\m$: $20=5t^2$, $t=2\s$; total $1+2=3\s$. \;(c) $v=gt=10(2)=20\ms$ (or $v^2=10^2+2(10)(15)=400$). \;(d) A straight line of slope $-10\mss$, from $+10\ms$ at $t=0$, through $0$ at $t=1\s$, to $-20\ms$ at $t=3\s$.
\end{solution}
\end{exo}
\Needspace{16\baselineskip}
\begin{exo}[Free response \textbullet\ Experimental design and analysis]{}
Students want to find the acceleration of a cart rolling down a ramp. They have a meter stick and a stopwatch.
\begin{enumerate}[label=(\alph*)]
\item Describe a procedure they could follow. Say which quantities they measure and how they reduce uncertainty.
\item They release the cart from rest and record the time $t$ to travel each distance $d$:
\begin{center}\begin{tabular}{c|cccc}
$d$ (m) & $0.25$ & $0.50$ & $0.75$ & $1.00$\\\hline
$t$ (s) & $0.50$ & $0.71$ & $0.87$ & $1.00$
\end{tabular}\end{center}
Which quantities should they plot to get a straight line? Explain using a kinematic equation.
\item Use the data to find the acceleration of the cart.
\end{enumerate}
\work{5cm}
\begin{solution}
(a) Mark several distances on the ramp; release the cart from rest at the top (no push); time the trip to each mark; repeat each measurement several times and average. \;(b) From rest, $d=\frac12at^2$: plot $d$ against $t^2$; the slope is $a/2$. \;(c) $t^2\approx0.25,\,0.50,\,0.76,\,1.00\,\mathrm{s^2}$; slope $\approx1.0\,\mathrm{m/s^2}$, so $a\approx2.0\mss$.
\end{solution}
\end{exo}
